\chapter{绪论}

\section{研究背景与研究意义}

固定翼飞行器着陆阶段是飞行全过程中风险相对较高的环节之一。与巡航阶段相比，着陆时飞行高度较低、速度变化较快、操纵时间较短，容易受到侧风、低空湍流、地面效应以及姿态扰动等因素的影响，从而导致接地点偏差、下沉率过大、姿态超限和滑跑偏移等问题。若控制不当，不仅会影响着陆质量，还可能对飞行器结构安全和后续运行带来不利影响。

随着固定翼飞行器特别是无人机在巡检、测绘、侦察、运输等领域广泛运用，飞行任务所处的环境变得越来越复杂，这对着陆控制系统的安全性、稳定性和适应性提出了更为严苛的要求。当前现有的着陆控制研究主要聚焦于控制算法、轨迹跟踪和姿态调节等方面，虽然能够在一定程度上提高着陆精度，但对于操纵者在着陆过程中的信息获取、风险判断和操作辅助所给予的关注相对欠缺。

在人机协同日益强化的背景下，把人机交互强化思想引入固定翼飞行器着陆控制系统，借助状态提示、风险告警、操纵辅助和安全约束等方式，能够帮助操纵者更精准地把握飞行状态，及时修正操作，从而降低着陆风险。因此，对基于人机交互强化的固定翼飞行器着陆控制系统进行研究，具有显著的现实意义。

从理论意义来看，本研究能够推动着陆控制研究从单一控制律设计向"控制层—交互层—保护层"协同设计方向延伸，丰富固定翼飞行器着陆控制系统的研究内容。从工程意义来看，本研究可为固定翼飞行器在复杂工况下实现安全着陆提供更具操作性的系统方案，为后续工程应用提供参考。

\section{国内外研究现状}

国外对于固定翼飞行器着陆的研究重点放在起落架结构设计、低速着陆气动特性、缓冲系统动力学和控制辅助等领域。Elahi等人针对无人作战飞行器CFRP起落架展开设计与分析，认为复合材料起落架能在保证承载能力的情况下实现轻量化。Lampropoulos等人针对低速无人机展开研究，发现低空区域存在的湍流和低速流场变化会对飞行器的升阻特性和姿态稳定性产生明显影响。Dinc等人对舰载无人机油气式起落架系统的动力学特性做了分析，探究了气体压缩、液体节流和结构非线性耦合作用对缓冲响应与冲击载荷的影响。Baek等人提出了基于输入信号变换的控制辅助系统，为把人机交互强化思想引入着陆控制系统提供了启发。

在国内，针对固定翼飞行器和无人机着陆的研究重点聚焦于起落架结构优化、整体布局分析和缓冲器性能研究等方面。帅涛针对某型号起落架进行了设计和模态分析，指出起落架的固有频率和振型对着陆冲击响应有重要影响。周著学和于萍借助有限元分析方法对小型航拍无人机起落架进行了优化设计。薛影、张超和吴星星剖析了无人机起落架布置对起飞性能的影响。罗杰等人研究了某型无人机起落架放下阶段缓冲器的气液流动特性。蒲雅熙从油气缓冲器气液两相阻尼特性着手，分析了不同工况下缓冲器阻尼变化规律。

目前国内外研究已为固定翼飞行器着陆控制奠定了良好基础，但现有研究大多是围绕结构、缓冲或控制中某一个方面进行的，针对"人机交互强化"的系统研究比较欠缺。因此本文在已有研究基础上引入人机交互强化思想，探究固定翼飞行器着陆控制系统的优化设计。

\section{研究内容与技术路线}

\subsection{研究内容}

本文以固定翼飞行器着陆过程为研究对象，围绕人机交互强化的着陆控制系统开展研究。首先，对着陆任务场景、系统组成及关键影响因素进行分析。其次，建立着陆阶段动力学模型和仿真模型。再次，设计模型校核方法、性能评价指标和典型工况，通过仿真计算分析不同工况下系统响应变化规律。最后，结合仿真结果提出人机交互强化策略。

\subsection{技术路线}

本文围绕"需求分析—模型建立—仿真验证—结果分析—策略优化"的思路进行研究。首先查阅相关文献，剖析着陆阶段主要风险和控制需求。接着建立动力学模型与仿真模型，完成参数设定。然后设计不同工况下的仿真方案，进行计算与比较分析。最后识别系统存在的问题，提出人机交互强化策略。
